% Generated draft: requires mathematical and source review.
\documentclass[11pt]{article}
\usepackage[margin=25mm]{geometry}
\usepackage{fontspec}
\setmainfont{Latin Modern Roman}
\usepackage{amsmath,amssymb,mathtools,unicode-math}
\setmathfont{Latin Modern Math}
\usepackage{enumitem,needspace,xurl,hyperref}
\setlength{\parindent}{0pt}
\setlength{\parskip}{5pt}
\setlength{\emergencystretch}{3em}
\newcommand{\N}{\mathbb{N}}
\newcommand{\Z}{\mathbb{Z}}
\newcommand{\Q}{\mathbb{Q}}
\newcommand{\R}{\mathbb{R}}
\newcommand{\C}{\mathbb{C}}
\newcommand{\F}{\mathbb{F}}
\newcommand{\A}{\mathbb{A}}
\newcommand{\PP}{\mathbb{P}}
\newcommand{\E}{\mathbb{E}}
\newcommand{\eps}{\varepsilon}
\newcommand{\dd}{\,\mathrm{d}}
\begin{document}
\section*{UnsolvedMath ID: 11100048}
\textbf{Source problem code:} catalogue entry
\subsection{Definitions and mathematical statement}
The source asks for a homotopical framework for sheaves over manifolds and its relation to de Rham cohomology, Floer homology, and quantum cohomology. A model structure here means specified classes of weak equivalences, cofibrations, and fibrations on a complete and cocomplete category, satisfying the Quillen model-category axioms. Specifying such a structure alone does not specify a cohomology theory: the relevant objects and homotopy-invariant or derived construction must also be given.

For the de Rham part, if $M$ is a smooth manifold, let $\Omega_M^{\bullet}$ be its complex of sheaves of real smooth differential forms. Its ordinary de Rham cohomology is
\[
 H^k_{\mathrm{dR}}(M)=H^k\bigl(\Gamma(M,\Omega_M^{\bullet}),d\bigr).
\]
The request is to identify an appropriate category of sheaves or sheaf complexes on $M$, a model structure on it, and a derived construction recovering this cohomology. The record does not fix sheaves of sets, spaces, modules, or differential graded objects, so the displayed de Rham complex specifies the familiar target without selecting the entire intended model-category framework.

For the Floer and quantum parts, the source asks whether the corresponding homology or cohomology theories can likewise be obtained from model structures. An answer would need to specify the variant of Floer homology or quantum cohomology, its geometric and analytic hypotheses and coefficient system, the category and model structure, and a comparison identifying its derived invariant with that theory. None of these data is fixed in the source. Consequently these are research-program requests rather than a single assertion about all manifolds or all Floer theories.
\subsection{Short English statement}
Find model-category explanations of de Rham cohomology and investigate analogous explanations of Floer homology and quantum cohomology. The supplied summary regards the de Rham component as understood through sheaf and differential graded model structures, while reporting the other components open as literally posed and pointing to differential graded or higher Fukaya categories as nearby frameworks. These are supplied assessments, with no precise construction or current verification provided.
\subsection{Sources}
\begin{itemize}
\item Local UnsolvedMath registry, permanent ID $11100048$, source code \texttt{catalogue entry}, proposed by Mark Hovey; created 2026-07-29, updated 2026-08-06. Status \texttt{partially\_solved}; research classification \texttt{PARTIAL-PROGRESS}.
\item Mark Hovey, \emph{Hovey's Algebraic Topology Problem List}, ``Model categories,'' item 2, surviving mirror \url{https://www-users.cse.umn.edu/~tlawson/hovey/}; accessed 2026-07-29, HTML ordered-list extraction. The background describes 65 top-level items across nine category pages and assigns default difficulty L3.
\item The original item motivates its questions through the analogy between a scheme and a ring, and between a manifold and $\R^n$, and expresses the author's source-era view that homotopy theory had made little impact on analysis partly through lack of effort. The summary mentions sheaf/DG model structures and the dg/$\infty$ Fukaya category, but supplies no exact works or authoritative theorem citations for its progress and nonexistence/open-status claims.
\item Source status evidence \texttt{NEEDS\_REVIEW}: originally presented as open, with current primary-literature status audit outstanding. Source rights \texttt{NEEDS\_REVIEW}: publicly accessible author/publisher copy, redistribution terms requiring review. Dataset attribution: Ulam AI / UnsolvedMath Contributors, CC BY 4.0.
\end{itemize}
\end{document}
